\section{Integer-exponent readings and bounded component multipliers}
\label{sec:sr-source-integer-reading}
This section follows the actual negative-exponent clearing argument,
rather than replacing it by an extension principle.
We distinguish the Fourier propositions used by the source's
bounded-section algebra from the algebraic section-reading argument.

\subsection{The bounded-section product}
Fix the source Hermitian frame datum \(X\), its open frame domain
\(\mathcal D_X\), and integral row weights \(\lambda,\lambda'\).
Write
\[
 \|P\|_{\lambda,x}=|P(x)|\,n_\lambda(x),\qquad
 n_\lambda(x)>0,\qquad
 n_{\lambda+\lambda'}(x)=n_\lambda(x)n_{\lambda'}(x).
 \tag{IR1}\label{eq:sr-source-reading-nfac}
\]
The last equality follows directly from the product of positive
successive Gram-minor ratios in \(n_\lambda\), using
\(r^{a+b}=r^ar^b\) for \(r>0\).
A bounded section means a holomorphic section with the source
Borel and arithmetic transformation laws and bounded weighted norm.
The source space \(SS_\lambda\) consists of such sections whose
zero Fourier coefficient vanishes along every admissible rational
isotropic direction.

Here are the two exact Fourier inputs used at this point:
\begin{enumerate}
\item A nonzero Fourier coefficient of a moderate holomorphic section
along an admissible isotropic direction has frequency
\(\nu=(\nu_\tau)\) with \(\nu_\tau\ge0\) for every \(\tau\).
A bounded section is moderate with exponent \(0\).
\item The coefficient of a product is the convergent convolution of
the two coefficient families:
\[
 \operatorname{fj}_{\lambda+\lambda'}(PQ;\xi,\nu,x)
 =\sum_{\nu'}\operatorname{fj}_\lambda(P;\xi,\nu',x)
       \operatorname{fj}_{\lambda'}(Q;\xi,\nu-\nu',x).
 \tag{IR2}\label{eq:sr-source-reading-convolution}
\]
The sum is the source's \(\mathrm{HasSum}\) assertion, not a
formal multiplication of series.
\end{enumerate}
These are exactly \srcid{ModerateCoefficientPositivity} and
\srcid{CoefficientConvolution}.

\begin{srlemma}\label{lem:sr-source-bdd-cusp-product}
Under these two Fourier inputs, \(P\in SS_\lambda\) and a bounded
section \(Q\) of weight \(\lambda'\) imply \(PQ\in SS_{\lambda+\lambda'}\).
\end{srlemma}
\begin{proof}
Holomorphy is preserved by multiplication.
Both transformation characters multiply; for the Borel character
this is the identity
\(\prod_{\tau,i}b_{\tau,ii}^{\lambda_{\tau,i}+\lambda'_{\tau,i}}
 =b_\lambda(b)b_{\lambda'}(b)\).
Every diagonal entry is nonzero.
Equation~\eqref{eq:sr-source-reading-nfac} gives
\(\|PQ\|_{\lambda+\lambda',x}
 =\|P\|_{\lambda,x}\|Q\|_{\lambda',x}\), so the product is bounded.

Apply \eqref{eq:sr-source-reading-convolution} at \(\nu=0\).
For a nonzero convolution term, the first input gives
\(\nu'_\tau\ge0\) and \(-\nu'_\tau\ge0\) for every \(\tau\).
Thus \(\nu'=0\), and the \(P\)-factor is zero by cuspidality.
Every term is therefore zero.
Uniqueness of the stated sum proves that the product's zero
coefficient vanishes. This is
\srcid{BoundedSectionAlgebra.siegelSpaceAlgebra_of}, at its
\(SS\)-multiplication clause.
\end{proof}
The same argument with the triangle inequality proves that \(SS\)
is closed under zero, addition and scalar multiplication.
In taking sums of coefficients, integrability on the finite-volume
period domain follows from boundedness: real isotropic flow
preserves all the Gram minors, hence \(n_\lambda\), and the
Fourier exponential has modulus one.
Consequently the source's span-defined \(\mathrm{holSub}\) equals
the space of families whose every component belongs to \(SS\).
This is \srcid{mem_holSub_model}.

\subsection{The actual integer-exponent reading}
Use the source \(E,H,D_w,l,X_o,c\) and an actual good complex
uniformization \(U\).
Suppose \(p+q\ge1\), a source key section
\(\kappa\in\mathrm{KeySec}(X_o,c)\) has been supplied, and
\(G=(G_1,G_2)\) has even \(G_2\) and antitone integer exponent
rows:
\[
 i\le j<\operatorname{mult}_{D_w}(\tau)
 \ \Longrightarrow\ G_1(\tau,j)\le G_1(\tau,i).
 \tag{IR3}\label{eq:sr-source-integer-antitone}
\]
Let \(f\) be an actual algebraic section in
\(\mathrm{secZ}(X_o,c,\mathrm{exZ}(D_w,G))\).
Retain the uniformization's actual reading law for dominant
natural weights and even tensor degree.
Also retain the actual nonvanishing section supplier:
for every \(x_0\in\mathcal D_{U,k}\) there are \(m\ge0\) and
\[
 f_0\in\mathrm{sec}(X_o,c,0,2(m+1)),\qquad
             \operatorname{rd}_U(f_0)_k(x_0)\ne0.
 \tag{IR4}\label{eq:sr-source-reading-nonvanishing}
\]
The source supplies (IR4) by its level-\(3\) eta rule in
\srcid{UnfoldedIsSecZeroKeySection.exists_sec_rdU_ne_zero}.

\begin{srproposition}\label{prop:sr-source-integer-reading}
For every component \(k\), the actual reading
\(\operatorname{rd}_U(f)_k\) is a section of the exact model weight
\(\operatorname{modelWt}(U,G,k)\).
\end{srproposition}
\begin{proof}
Fix \(x_0\) and choose (IR4).
Write \(G_2=2g_2\), set \(N=|G_2|\), and choose \(B\ge0\) with
\(-B\le G_1(\tau,j)\) for every actual row index.
The finite sum of all \(|G_1(\tau,j)|\) is one possible \(B\).
Set exactly as in the source
\[
 \begin{aligned}
 t_g&=(B+1)m+(N+1)(m+1),&t&=g_2+t_g,\\
 a_g(\tau,i)&=
 \begin{cases}B+1&i<\operatorname{mult}_{D_w}(\tau),\\0&\text{otherwise},
 \end{cases}&
 a(\tau,i)&=
 \begin{cases}G_1(\tau,i)+B+1&i<\operatorname{mult}_{D_w}(\tau),\\
 0&\text{otherwise}.
 \end{cases}
 \end{aligned}
 \tag{IR5}\label{eq:sr-source-integer-clearing}
\]
The lower bound on \(G_1\) makes \(a\) natural-valued.
Moreover \(t_g\ge N+1\), so \(t\ge0\).
Both \(a_g\) and \(a\) are antitone, including the transition
to the zero entries outside the actual rows.

Let \(\rho=\kappa.\rho\), the specified algebraic unit of key
exponent, and put
\[
                  g=\rho^{B+1}f_0^{B+N+2}.
 \tag{IR6}\label{eq:sr-source-reading-clearing-section}
\]
The exact exponent identities are
\[
 \begin{split}
 \mathrm{exZ}(D_w,\mathrm{wtZ}(a_g,2t_g))
  &=(B+1)\,\mathrm{keyEx}(D_w)
      +(B+N+2)\mathrm{exZ}(D_w,\mathrm{wtZ}(0,2(m+1))),\\
 \mathrm{exZ}(D_w,\mathrm{wtZ}(a,2t))
  &=\mathrm{exZ}(D_w,G)
       +\mathrm{exZ}(D_w,\mathrm{wtZ}(a_g,2t_g)).
 \end{split}
 \tag{IR7}\label{eq:sr-source-reading-exponent-identities}
\]
For the tensor coordinate the first identity is
\((B+1)(-2)+2(B+N+2)(m+1)=2t_g\);
the second uses \(G_2=2g_2\).
For every determinant coordinate they read respectively
\(B+1=B+1\) and
\(G_1+B+1=G_1+(B+1)\).
Multiplication and powers of actual algebraic sections therefore give
\[
       g\in\mathrm{sec}(a_g,2t_g),\qquad
       fg\in\mathrm{sec}(a,2t).
\]
The uniformization's given dominant-weight reading law makes their
readings sections of the corresponding natural model weights.

The model weights add, because both original tensor weights are
even and their algebraic exponents agree in (IR7).
Thus, writing \(\mu=\operatorname{modelWt}(U,G,k)\) and
\(\lambda=\mathrm{laU}(U,a_g,2t_g,k)\), the readings of
\(fg\) and \(g\) have weights \(\mu+\lambda\) and \(\lambda\).

Their denominator reading at \(x_0\) is nonzero.
Indeed \(\operatorname{rd}_U(f_0)_k(x_0)\ne0\).
Multiplicativity then forces the reading of \(1\) to be \(1\)
there. Since \(\rho\rho^{-1}=1\), the reading of \(\rho\) is
nonzero there too. Equation (IR6) now gives the assertion.
Finally,
\[
    \operatorname{rd}_U(f)_k\,\operatorname{rd}_U(g)_k
                     =\operatorname{rd}_U(fg)_k
\]
holds on the whole frame domain.

For completeness, the source's local quotient argument gives
all three section clauses. Near \(x_0\) the denominator stays
nonzero, so the quotient is holomorphic.
For an upper triangular invertible \(b\),
\(b_\lambda(b)\ne0\).
Evaluate the displayed identity at \(x_0\) and \(x_0b\) and
cancel \(b_\lambda(b)\operatorname{rd}_U(g)_k(x_0)\);
this gives the Borel weight \(\mu\).
For an arithmetic transformation \(\gamma\), the multiplier
\(\chi_\lambda(\gamma)\) is nonzero, since it has finite order.
The identical cancellation gives the multiplier
\(\chi_\mu(\gamma)\).
Since \(x_0\) was arbitrary, all section clauses hold everywhere.
These are precisely
\srcid{IsSecReadingDecreasingIntegerExponents.isSec_of_local_quot}
and \srcid{isSec_modelWt_rdU}.
\end{proof}

\subsection{Preservation of the actual component cusp space}
Let \(Z\) be the source component, and let \(s\) have the same
hypotheses as \(f\) in the preceding proposition, with weight \(G_0\).
Assume also
\[
 s\in\mathrm{suppOn}(Z),\qquad
 \forall\sigma\text{ good}\ \exists C_\sigma\
 \forall z\in\mathrm{framesOn}(U_\sigma,Z):\
                    \mathrm{homN}(s;z)\le C_\sigma .
 \tag{IR8}\label{eq:sr-source-component-bdd}
\]
All frames and norms here belong to the same supplied
uniformization. No uniformity in a varying family is being added.

\begin{srproposition}\label{prop:sr-source-component-cusp-multiplier}
Under the two Fourier inputs above, for every \(G\) with even
tensor coordinate,
\[
      f\in\mathrm{cuspOn}(U,Z,G)
          \quad\Longrightarrow\quad
      fs\in\mathrm{cuspOn}(U,Z,G+G_0).
 \tag{IR9}\label{eq:sr-source-component-cusp-product}
\]
\end{srproposition}
\begin{proof}
Algebraic section exponents add, and support of \(s\) on \(Z\)
gives support of \(fs\) there.
Fix a good \(\sigma\).
The actual model reading of \(s\) is a section by the preceding
proposition.
On \(\mathrm{framesOn}(U_\sigma,Z)\), its weighted norm is at most
\(C_\sigma/\mathrm{lieFacZ}(D_w,x,G_0)\), whose denominator is
strictly positive.
Off those component frames its reading is zero by the support
condition. It is therefore a bounded section on every component,
with bound
\[
                \max\{0,C_\sigma/\mathrm{lieFacZ}(D_w,x,G_0)\}.
\]
The component readings of \(f\) belong to \(SS\).
The first lemma puts their products in \(SS\) at the sum model
weight. The equality between \(\mathrm{holSub}\) and componentwise
\(SS\), followed by the source definition of
\(\mathrm{cuspW}\) and the already checked support, proves (IR9).
This is exactly
\srcid{ExistsUnramPtsDefinite.mul_mem_cuspOn_of_isSec}
followed by
\srcid{MulMemCuspFormsOfBoundedSectionComponent.mul_mem_cuspOn_of_bdd}.
\end{proof}

\paragraph{Frozen-source ledger and retained suppliers.}
All paths below use \srcid{alpoge-source-7cbf7a0c}.
In
\srcid{src/Development/AbelianVarietyTheory/Proofs/DivisorProofs/},
the actual reading proof is
\srcid{IsSecReadingDecreasingIntegerExponents.lean},
blob \srcid{b08455118b8cd54243ae8185790eaef6bf4e9382};
the final bounded component assembly is
\srcid{MulMemCuspFormsOfBoundedSectionComponent.lean},
blob \srcid{3a5a50359e5689723974601a420411cab829aa03}.
The earlier componentwise bounded-product argument is in
\srcid{src/Development/AbelianVarietyTheory/Definitions/Divisors/ExistsUnramPtsDefinite.lean},
blob \srcid{3cd59129266bb5e4fbc6a994794b40622e259d41}.
The actual nonvanishing supplier is
\srcid{src/Development/AbelianVarietyTheory/Proofs/IwasawaDecomposition/UnfoldedIsSecZeroKeySection.lean},
blob \srcid{54d610b59fda59d5c0fb5d17e846651747cf1cd4}.
Its level-\(3\) eta construction and the existence of the key unit
are separate actual inputs to the reading proposition.

The Fourier product proof is in
\srcid{src/ReciprocityTheory/ReductiveGroups/ArithmeticGeometry/FundamentalDomains/BoundedSectionAlgebra.lean},
blob \srcid{71f99c21a0d1bc7a47618506aea05641e5cad761}.
The concrete Fourier suppliers remain
\srcid{ModerateCoefficientPositivityHolds.lean}
(blob \srcid{712132fb6d8bf919371842e7a2a95d5f42d1fc57})
and
\srcid{CoefficientConvolutionHolds.lean}
(blob \srcid{dbd4d7e005c93fdc46adbd27159fb35b54cc2fbb})
in that same directory.
The latter itself uses
\srcid{CoefficientConvolutionOfSubCoefficients}
and \srcid{SublatticeCoefficients}.
This section proves the actual exponent clearing, local quotient,
and bounded component multiplication from those precise suppliers;
it does not substitute an unnamed Koecher principle for them.
